\subsection{扩散}

定义 3. --- 设 X 是一个局部紧空间，\(\Lambda : t \mapsto \lambda_t\) 是 T 到 \(\mathcal{M}_{+}(X)\) 中的一个映射。称映射 \(\Lambda\) 是 T 在 X 中的一个扩散，如果 \(\Lambda\) 对 T 上每个具有紧支集的正测度都是适宜的。称扩散 \(\Lambda\) 是有界的，如果所有测度 \(\lambda_{t}\) 都有界且 \(\sup_{t\in T}\|\lambda_{t}\|<+\infty\) ；这个量称为 \(\Lambda\) 的范数，记作 \(\|\Lambda\|\) 。

下面的命题只是定义的转述：

命题 8. --- 映射 \(\Lambda: t \mapsto \lambda_t\) （T 到 \(\mathcal{M}_+(\mathrm{X})\) 中的）是一个扩散，必须且只须下列条件成立：

\begin{enumerate}
\def\labelenumi{\arabic{enumi})}
\item
  对每个下半连续函数 \(f \geqslant 0\) （定义在 \(X\) 上），函数 \(t \mapsto \lambda_t^\bullet(f)\) 在 \(T\) 上是普遍可测的。
\item
  对每个函数 \(g \in \mathcal{K}_{+}(X)\) ，函数 \(t \mapsto \lambda_t(g)\) 在 \(T\) 中是局部有界的。
\item
  对每个定义在 X 上的下半连续函数 \(f \geqslant 0\) 以及 T 上每个具有紧支集的正测度 \(\mu\) ，下列关系式成立，其中 \(\nu\) 表示 \(\int \lambda_{t} d\mu(t)\) ：
\end{enumerate}

\[
\int^ {\bullet} f (x) d \nu (x) = \int^ {\bullet} d \mu (t) \int^ {\bullet} f (x) d \lambda_ {t} (x).\tag{12}
\]

假设 \(\Lambda\) 是一个扩散。则由适宜映射的定义（No.~1，定义 1）及命题 6，条件 1) 成立；由公式 (4)，条件 3) 成立，因为 \(\Lambda\) 是 \(\mu\) -适宜的。设 \(g \in \mathcal{K}_{+}(X)\) ，\(u\) 是函数 \(t \mapsto \lambda_t(g)\) （由 1)，它是普遍可测的）；假设 \(u\) 不是局部有界的。那么将存在一个紧集 \(K\) 使 \(u\) 在 \(K\) 上不有界，从而将存在元素的一个序列 \((t_n)\) （取自 \(K\) ），使 \(u(t_n) \geqslant n^2\) 对所有 \(n \geqslant 1\) 成立；于是 \(u\) 对具有紧支集的测度 \(\mu = \sum_{n \geqslant 1} \frac{1}{n^2} \varepsilon_{t_n}\) 不是可积的，这与关于 \(\Lambda\) 的假设相矛盾，该假设蕴涵 \(t \mapsto \lambda_t(g)\) 对每个具有紧支集的正测度可积。因此上述三个条件是必要的。反之，条件 1) 与 2) 蕴涵 \(\Lambda\) 对每个具有紧支集的测度 \(\mu\) 都是标量本质 \(\mu\) -可积的。由于每个测度 \(\mu' \geqslant 0\) （被一个具有紧支集的测度 \(\mu\) 优超者）本身也具有紧支集，条件 1) 与 3) 表明 \(\Lambda\) 对每个具有紧支集的正测度都是 \(\mu\) -适宜的，这正是所要的结果。

命题 9. --- 设 \(\Lambda : t \mapsto \lambda_t\) 是 T 到 \(\mathcal{M}_+(\mathrm{X})\) 中的一个映射，使得函数 \(t \mapsto \lambda_t(g)\) 对每个 \(g \in \mathcal{K}_+(\mathrm{X})\) 都在 T 上普遍可测且局部有界。在下列每一种情形，都可以断言 \(\Lambda\) 是一个扩散：

\begin{enumerate}
\def\labelenumi{\alph{enumi})}
\item
  X 的拓扑具有可数基；
\item
  \(\Lambda\) 关于淡拓扑是普遍可测的。
\end{enumerate}

因为，设 \(\mu\) 是 T 上一个具有紧支集的正测度；映射 \(\Lambda\) 是标量本质 \(\mu\) -可积的，因而只要 a) 或 b) 之一成立，它就是 \(\mu\) -适宜的（No.~1，命题 2）。

在本节其余部分，我们将采用以下记号：我们将用 \(\langle\eta,h\rangle\) 表示正测度 \(\eta\) 下的正 \(\eta\) -可测函数 h 的本质上积分。映射 \(\Lambda:t\mapsto\lambda_{t}\) 将是 T 在 X 中的一个扩散。如果 f 是定义在 X 上的一个正的普遍可测函数，我们将用 \(\Lambda f\) 表示映射 \(t\mapsto\lambda_{t}^{\bullet}(f)\) 。如果 \(\mu\) 是 T 上使 \(\Lambda\) 为标量本质 \(\mu\) -可积的正测度，我们将用 \(\mu\Lambda\) 表示测度 \(\int\lambda_{t}d\mu(t)\) 。于是积分的定义取如下形式

\[
\langle \mu \Lambda , f \rangle = \langle \mu , \Lambda f \rangle \quad \text { for } f \in \mathscr {K} _ {+} (\mathrm{X}).
\]

我们将说 T 上的正测度 \(\mu\) 属于 \(\Lambda\) 的定义域，如果 \(\Lambda\) 是 \(\mu\) -适宜的：这等价于说（依据命题 8）\(\Lambda\) 是标量本质 \(\mu\) -可积的，且 \(\langle\mu'\Lambda,f\rangle=\langle\mu',\Lambda f\rangle\) 对每个正测度 \(\mu'\leqslant\mu\) 及每个正的下半连续函数 f 成立。

命题 10. --- 设 f, g 是 X 上两个正的普遍可测函数，a 是一个数 \(\geqslant 0\) ，\(\mu\) 和 \(\nu\) 是 T 上两个正测度。那么：

\begin{enumerate}
\def\labelenumi{\alph{enumi})}
\item
  \(\Lambda (f + g) = \Lambda f + \Lambda g, \Lambda (af) = a\Lambda f.\)
\item
  如果 \(\mu\) 和 \(\nu\) 属于 \(\Lambda\) 的定义域，那么 \(\mu + \nu\) 和 \(a\mu\) 也属于其定义域，并且有 \((\mu + \nu)\Lambda = \mu \Lambda + \nu \Lambda\) ，\((a\mu)\Lambda = a(\mu \Lambda)\) 。
\end{enumerate}

唯一不明显之点是 \(\mu + \nu\) 属于 \(\Lambda\) 的定义域，这一点可以这样处理：注意到每个被 \(\mu + \nu\) 优超的正测度都具有 \(\mu' + \nu'\) 的形式，其中 \(\mu' \leqslant \mu\) ，\(\nu' \leqslant \nu\) （“分解引理”，Ch. II, §1, No.~1）。另见下一命题。

命题 11. --- 正测度 \(\mu\) （\(T\) 上的）属于 \(\Lambda\) 的定义域，必须且只须 \(\Lambda\) 是标量本质 \(\mu\) -可积的。

该条件显然是必要的。反之，设它成立，并设 f 是定义在 X 上的一个正的下半连续函数。函数 \(\Lambda f\) 是普遍可测的，因而是 \(\mu\) -可测的。我们将证明 \(\langle\mu,\Lambda f\rangle=\langle\mu\Lambda,f\rangle\) ；由于这个等式对每个正测度 \(\mu'\leqslant\mu\) 也成立（因为 \(\Lambda\) 也是标量本质 \(\mu'\) -可积的），便可推出 \(\Lambda\) 是 \(\mu\) -适宜的。

设 \((\mu_i)_{i\in I}\) 是一个由具有紧支集的正测度组成的可和族，使得 \(\mu = \sum_{i\in I}\mu_i\) （§2, No.~3, 命题 4）；于是测度族 \(\mu_{i}\Lambda\) 是可和的，且 \(\mu \Lambda = \sum_{i\in I}\mu_{i}\Lambda\) （No.~1, 命题 1 的推论）。从而 \(\langle \mu \Lambda ,f\rangle = \sum_{i\in I}\langle \mu_i\Lambda ,f\rangle\) （§2, No.~2, 命题 1）；但 \(\Lambda\) 是 \(\mu_{i}\) -适宜的，故 \(\langle \mu_i\Lambda ,f\rangle = \langle \mu_i,\Lambda f\rangle\) 。再应用 §2 的命题 1，就得到所要的等式：

\[
\langle \mu \Lambda , f \rangle = \sum_ {i \in I} \langle \mu_ {i} \Lambda , f \rangle = \sum_ {i \in I} \langle \mu_ {i}, \Lambda f \rangle = \langle \mu , \Lambda f \rangle .
\]

推论 1. --- 如果 \(\Lambda\) 是一个有界扩散，那么每个有界正测度 \(\mu\) 都属于 \(\Lambda\) 的定义域，且 \(\|\mu\Lambda\| \leqslant \|\mu\|\|\Lambda\|\) 。

推论 2. --- 假设 \(\mu\) 是一个可和族 \((\mu_{\alpha})_{\alpha \in \mathbb{A}}\) 之和，该族由属于 \(\Lambda\) 的定义域的正测度组成。\(\mu\) 属于 \(\Lambda\) 的定义域，必须且只须测度族 \(\mu_{\alpha}\Lambda\) 可和，此时 \(\mu \Lambda = \sum \mu_{\alpha}\Lambda\) 。

只需应用 No.~1 的命题 1 的推论。

命题 5 用扩散的语言表述，就取如下的形式：

命题 12. --- 设 \(\mu\) 是 T 上属于 \(\Lambda\) 的定义域的一个正测度，\(f\) 是定义在 X 上的一个普遍可测函数 \(\geqslant 0\) 。如果 \(f\) 关于测度 \(\mu\Lambda\) 是适度的，或者如果诸测度 \(\lambda_t\) 有界，那么函数 \(\Lambda f\) 是 \(\mu\) -可测的，并且

\[
\langle \mu \Lambda , f \rangle = \langle \mu , \Lambda f \rangle .\tag{13}
\]

推论. --- 如果 X 在无穷远处可数，或者如果诸测度 \(\lambda_{t}\) 有界，那么函数 \(\Lambda f\) 对每个定义在 X 上的普遍可测函数 \(f \geqslant 0\) 都在 T 上普遍可测，且 (13) 成立。

